\documentclass[UTF8,a4paper,12pt]{ctexart}

\usepackage{geometry}
\usepackage{graphicx}
\usepackage{hyperref}
\usepackage{enumitem}
\usepackage{xcolor}

\geometry{left=2.5cm,right=2.5cm,top=2.5cm,bottom=2.5cm}
\hypersetup{
  colorlinks=true,
  linkcolor=blue,
  urlcolor=blue
}
\setlist[itemize]{leftmargin=2em}

\title{算法基础第一次大作业报告}
\author{你的姓名}
\date{2026年5月}

\begin{document}

\maketitle

\section{个人主页介绍}

\subsection{主页访问链接}

主页访问链接：\url{https://example.com}

\subsection{主页主要内容}

\begin{itemize}
  \item 个人基本介绍：这里填写你的真实姓名、专业背景、年级或学习状态。
  \item 学习方向或技术兴趣：这里填写你目前关注的方向，例如算法、数据结构、AI Agent、三维视觉、计算机图形学等。
  \item 项目展示或学习记录：这里填写课程作业、项目、阅读记录或阶段性学习总结。
  \item 博客入口：主页中提供了第一篇技术博客的入口，能够直接跳转到博客页面。
\end{itemize}

\subsection{页面结构与功能说明}

个人主页采用静态网页实现，主要包括顶部导航、个人简介、学习方向、项目或学习记录、博客入口等模块。页面链接可以正常跳转，并考虑了基本的响应式布局，便于在常见浏览器和不同屏幕尺寸下阅读。

\section{博客内容介绍}

\subsection{博客主题}

博客主题：我如何理解 AI Agent 辅助编程

\subsection{博客访问方式}

博客访问链接：\url{https://example.com/blog/first-blog.html}

\subsection{主要内容概述}

博客主要讨论 AI Agent 辅助编程与传统网页端问答的区别，并结合本次个人主页搭建过程，说明为什么需要项目规范文件、版本管理、验证清单和交接记录。

\subsection{个人理解与思考}

这里填写你自己的理解。建议结合本次作业中的具体经历，例如：AI 帮助你完成了哪些页面结构或样式设计，你自己如何检查、修改和取舍 AI 生成的内容。

\section{实现过程}

\subsection{使用的主要工具或技术}

\begin{itemize}
  \item AI Agent：Codex。
  \item 技术栈：HTML、CSS，必要时使用少量 JavaScript。
  \item 版本管理：Git。
  \item 部署方式：GitHub Pages。
\end{itemize}

\subsection{网页搭建过程}

首先阅读课程大作业要求，确定需要完成个人主页、技术博客和作业报告。随后建立 Git 仓库，编写 \texttt{AGENTS.md} 作为 AI Agent 的项目规范文件，明确项目目标、文件结构、内容约束、验证清单和交接格式。之后逐步实现首页、博客页面和样式文件。

\subsection{部署过程}

这里填写部署到 GitHub Pages 的过程。正式部署后，需要补充 GitHub 仓库地址、Pages 设置方式和最终访问链接。

\subsection{遇到的问题及解决方法}

这里填写实现过程中遇到的问题。例如：路径跳转错误、移动端布局不合适、Git 初始化或部署配置问题，以及对应的解决方法。

\section{AI Agent 使用说明}

\subsection{使用了哪些 AI 工具}

本次作业主要使用 Codex 作为 AI Agent，辅助阅读作业要求、整理项目结构、编写初始网页和报告模板。

\subsection{AI 辅助完成了哪些内容}

\begin{itemize}
  \item 提取并总结课程 PDF 中的大作业要求。
  \item 结合助教讲义整理 AI Agent 工作流。
  \item 生成并修改 \texttt{AGENTS.md} 项目规范文件。
  \item 搭建静态个人主页和第一篇博客页面的初始版本。
  \item 建立报告模板和 AI 工作记录文档。
\end{itemize}

\subsection{自己对 AI 生成内容做了哪些修改}

这里填写你自己完成的修改，例如：替换个人真实信息、调整页面内容、选择博客主题、补充个人思考、检查链接和截图。

\subsection{对 AI 辅助编程的思考}

这里填写你的反思。建议强调：AI 可以提高搭建和整理效率，但最终内容真实性、技术判断、结果检查和提交质量仍然需要自己负责。

\section{结果展示}

\subsection{个人主页截图}

正式提交前，将主页截图放到 \texttt{report/images/} 目录，并在这里引用。

% 示例：
% \begin{figure}[htbp]
%   \centering
%   \includegraphics[width=0.9\textwidth]{images/homepage.png}
%   \caption{个人主页截图}
% \end{figure}

\subsection{博客页面截图}

正式提交前，将博客页面截图放到 \texttt{report/images/} 目录，并在这里引用。

% 示例：
% \begin{figure}[htbp]
%   \centering
%   \includegraphics[width=0.9\textwidth]{images/blog.png}
%   \caption{博客页面截图}
% \end{figure}

\section{提交信息}

\begin{itemize}
  \item 个人主页网址：\url{https://example.com}
  \item 博客入口：\url{https://example.com/blog/first-blog.html}
  \item 报告导出日期：2026年5月
\end{itemize}

\end{document}
